% concatenated from main.tex
 %\pdfoutput=1
\documentclass[12pt,a4paper,reqno]{amsart}
\usepackage{amssymb}
\usepackage{amscd}
%\usepackage[pdftex,pdfpagelabels]{hyperref}
\usepackage{enumerate}
%\usepackage{psfig}
\usepackage{graphicx}
%\usepackage{showkeys}  
\usepackage{siunitx}
\usepackage{tikz-cd}
\usepackage{color}
\usetikzlibrary{arrows}
% uncomment this when editing cross-references
\numberwithin{equation}{section}

\usepackage{mathabx}


\usepackage{mathtools}%                  http://www.ctan.org/pkg/mathtools
\usepackage[tableposition=top]{caption}% http://www.ctan.org/pkg/caption
\usepackage{booktabs,dcolumn}%           http://www.ctan.org/pkg/dcolumn + http://www.ctan.org/pkg/booktabs
\usepackage{hyperref}
% Lighter notation.
%\newcommand*\mc[1]{\multicolumn{1}{c}{#1}}
%\newcommand*\tupref[2]{\href{http://math.mit.edu/~primegaps/tuples/admissible_#1_#2.txt}{\num{#2}}}



%\DeclareMathOperator*\Kl{Kl} (commented because yield bad display for  \Kl_q %replaced with \newcommand... )
%\DeclareMathOperator*\FT{FT} (commented because yield bad display for \FT_q
%replaced with \newcommand...)

\DeclareMathOperator*\swan{swan}
\DeclareMathOperator*\cond{cond}
\DeclareMathOperator*\Gal{Gal}

\newcommand{\FT}{\mathrm{FT}}
\newcommand{\Kl}{\mathcal{K}\ell}
% Setup for ``caption''.
%\DeclareCaptionLabelSeparator{separation}{:\quad}
%\captionsetup{
	%font=small,
	%labelfont=sc,
	%labelsep=separation,
	%width=0.8\textwidth
	%}

\DeclareFontFamily{OT1}{rsfs}{}
\DeclareFontShape{OT1}{rsfs}{n}{it}{<-> rsfs10}{}
\DeclareMathAlphabet{\mathscr}{OT1}{rsfs}{n}{it}

\addtolength{\textwidth}{3 truecm}
\addtolength{\textheight}{1 truecm}
\setlength{\voffset}{-.6 truecm}
\setlength{\hoffset}{-1.3 truecm}

\theoremstyle{plain}

\newtheorem{theorem}{Theorem}[section]
%\newtheorem{theorem}[theorem]{Theorem}
\newtheorem{proposition}[theorem]{Proposition}
\newtheorem{lemma}[theorem]{Lemma}
\newtheorem{corollary}[theorem]{Corollary}
\newtheorem{conjecture}[theorem]{Conjecture}
\newtheorem{principle}[theorem]{Principle}
\newtheorem{claim}[theorem]{Claim}

\theoremstyle{definition}

%\newtheorem{roughdef}[subsection]{Rough Definition}
\newtheorem{definition}[theorem]{Definition}
\newtheorem{remark}[theorem]{Remark}
\newtheorem{remarks}[theorem]{Remarks}
\newtheorem{example}[theorem]{Example}
\newtheorem{examples}[theorem]{Examples}
%\newtheorem{problem}[subsection]{Problem}
%\newtheorem{question}[subsection]{Question}

\newcommand\F{\mathbb{F}}
\newcommand\E{\mathbb{E}}
\newcommand\R{\mathbb{R}}
\newcommand\Z{\mathbb{Z}}
\newcommand\N{\mathbb{N}}
\newcommand\n{\mathbf{n}}
\newcommand\one{\mathbf{1}}
\newcommand\C{\mathbb{C}}
\newcommand\Q{\mathbb{Q}}
\newcommand\eps{\varepsilon}
\newcommand\B{\operatorname{B}}
\newcommand\SO{{\operatorname{SO}}}
\newcommand\Rot{{\operatorname{Rot}}}
\newcommand\Dil{{\operatorname{Dil}}}
\newcommand\df{{\operatorname{df}}}
\newcommand\Dcal{\mathcal{D}}
%\newcommand\DHL{\operatorname{DHL}}
%\newcommand\EH{\operatorname{EH}}
%\newcommand\MPZ{\operatorname{MPZ}}
\newcommand\TypeI{\operatorname{Type}_{\operatorname{I}}}
\newcommand\TypeII{\operatorname{Type}_{\operatorname{II}}}
\newcommand\TypeIII{\operatorname{Type}_{\operatorname{III}}}
\renewcommand\P{\mathbb{P}}

\newcommand\rk{\mathrm{rk}}
\newcommand\Fr{\mathrm{Fr}}
\newcommand\tr{\mathrm{tr}}

\newcommand\ov{\overline}
\newcommand\wcheck{\widecheck}
%\newcommand\cond{\mathrm{cond}}
%\newcommand\Gal{\mathrm{Gal}}
%\newcommand\Kl{\mathrm{Kl}}
%\newcommand\swan{\mathrm{swan}}
%\newcommand\FT{\mathrm{FT}}
\newcommand\mcF{\mathcal{F}}
\newcommand\mcG{\mathcal{G}}
\newcommand\mcK{\mathcal{K}}
\newcommand\mcC{\mathcal{C}}
\newcommand\mcL{\mathcal{L}}
\newcommand\Gm{\mathbb{G}_m}
\newcommand\Af{\mathbb{A}^1}
\newcommand\Aa{\mathbb{A}}
\newcommand\Fp{\F_{p}}
\newcommand\Fq{\F_{q}}
\def\stacksum#1#2{{\stackrel{{\scriptstyle #1}}
		{{\scriptstyle #2}}}}
\newcommand{\mods}[1]{\,(\mathrm{mod}\,{#1})}
\newcommand{\piar}{\pi_1}
\newcommand{\pige}{\pi_1^{g}}
\newcommand\sumsum{\mathop{\sum\sum}\limits}
%\newcommand\E{\mathbb{E}}
\newcommand\Rr{\mathbb{R}}
\newcommand\Zz{\mathbb{Z}}
\newcommand\Nn{\mathbb{N}}
\newcommand\Cc{\mathbb{C}}
\newcommand\Qq{\mathbb{Q}}
\newcommand{\what}{\widehat}
\newcommand{\ra}{\rightarrow}

\newcommand{\alert}[1]{{\bf \color{red} TODO: #1}}

\newcommand{\on}{\operatorname}
\newcommand{\Log}{\on{Log}}

\parindent 0mm
\parskip   5mm 

%PhM use this for navigating the pdf when editing with texpad (comment if you prefer)
% (Ah, figured out why my LaTeX was not liking hyperref - there were old aux files that didn't use hyperref. -TT)
%\usepackage{hyperref}

% for optimize.tex
\usepackage{algorithm, algorithmic}

\usepackage[colorinlistoftodos]{todonotes}

\begin{document}
	
	\title{If $\sum_n n! c_n z^n$ is entire and $c_n$ does not terminate, then $\sum_n c_n z^n$ has infinitely many zeros}
	
	\author{Alann Rosas}
	\address{UC Irvine Department of Mathematics, Irvine, CA 92617.}
	\email{alannr@uci.edu}
	
	%\email{}
	
	\subjclass[2010]{30D20}
	
	\begin{abstract} We prove that if $\sum_n n! c_n z^n$ is entire and $c_n$ does not terminate, then $\sum_n c_n z^n$ has infinitely many zeros. We then use this result to give alternative proofs that the Le Roy functions $f_r(z)=\sum_{n=0}^\infty \frac{z^n}{(n!)^r}$ for $r>1$ and Bessel functions $J_\alpha(z)=\sum_{m=0}^\infty \frac{(-1)^m}{m!\Gamma(m+\alpha+1)}\left(\frac{z}{2}\right)^{2m+\alpha}$ for $\alpha\in\R$ have infinitely many zeros.
	\end{abstract}

    \keywords{Entire functions, complex analysis, infinitely many zeros, Maclaurin coefficients}
	
	\maketitle
	%\today
	
	
	%%%%%%%%%%%%%%%%%%%%%%%%%%%%%%%%%%%%%%%%%%%%%%%%%

    \section{Main Result}
    The following theorem is the main result of this paper:
    \begin{theorem}\label{thm1.1}
        Suppose $\sum_{n=0}^\infty n!c_n z^n$ is entire and the coefficients $c_n$ do not terminate. Then the function $f(z)=\sum_{n=0}^\infty c_n z^n$ is entire and has infinitely many zeros.
    \end{theorem}
    This gives a simple method for determining if a power series $\sum_n c_n z^n$ has infinitely many zeros. To illustrate its applicability, we give in Section \ref{sec3} alternative proofs that the Le Roy functions\footnote{For more information on the Le Roy functions, the reader is invited to have a look at \cite{le-roy}}
    $$f_r(z)=\sum_{n=0}^\infty \frac{z^n}{(n!)^r},\text{ }r>1$$
    and Bessel functions
    $$J_\alpha(z)=\sum_{m=0}^\infty \frac{(-1)^m}{m!\Gamma(m+\alpha+1)}\left(\frac{z}{2}\right)^{2m+\alpha}$$
    have infinitely many zeros.

    To prove Theorem \ref{thm1.1}, we will need two lemmas:
    \begin{lemma}\label{lem1.2}
        Let $f$ be an entire function of order $1$ with finitely many zeros. Let $N$ be the number of zeros counting multiplicity. Then there exists a constant $k\neq 0$ and a degree-$N$ polynomial $Q(n)$ (depending on $k$) such that for $n>N$, we have $f^{(n)}(0)=k^nQ(n)$.
    \end{lemma}
    \begin{lemma}\label{lem1.3}
        Suppose that $\sum_{n=0}^\infty n!c_n z^n$ converges on some neighborhood of $0$. Then $f(z)=\sum_{n=0}^\infty c_n z^n$ is entire and has order at most $1$.
    \end{lemma}
    Before proving these lemmas, let us show how they can be applied to give the proof of Theorem \ref{thm1.1}:

    \begin{proof}
        Since $\sum_{n=0}^\infty n!c_n z^n$ is entire and \textit{a fortiori} convergent on some neighborhood of $0$, we have by Lemma \ref{lem1.3} that $f$ is entire and has order at most $1$. Suppose, for the sake of reaching a contradiction, that $f$ has \textit{finitely many} zeros. Then the Hadamard factorization of $f$ is
        $$f(z)=e^{A(z)}P(z)$$
        where $A(z)$ and $P(z)$ are polynomials with $A(z)$ either constant or linear. Note that $A(z)$ cannot be constant, otherwise $f$ is a polynomial and hence $c_n=f^{(n)}(0)/n!$ terminates, contradicting our assumption. Thus, $A(z)=l+kz$ where $k\neq 0$, and hence
        $$f(z)=e^{l+kz}P(z)=e^{kz}\tilde{P}(z)$$
        where $\tilde{P}(z):=e^l P(z)$ is a polynomial. Let $N$ be the degree of $\tilde P$, so $N$ is also the number of zeros of $f$ counting multiplicity. By Lemma \ref{lem1.2}, there is a degree-$N$ polynomial $Q(n)$ such that $f^{(n)}(0)=k^n Q(n)$ holds for all $n>N$, so we can write
        $$\sum_{n=N+1}^\infty n!c_n z^n = \sum_{n=N+1}^\infty f^{(n)}(0) z^n=\sum_{n=N+1}^\infty Q(n) (kz)^n$$
        Noting that $Q(n)$ is a polynomial and thus $Q(n)\neq 0$ for all $n$ sufficiently large, we may determine the convergence radius of $\sum_{n=N+1}^\infty Q(n) (kz)^n$ using the ratio test:
        $$\lim_{n\to\infty}\left|\frac{Q(n+1) (kz)^{n+1}}{Q(n) (kz)^n}\right|=|kz|\lim_{n\to\infty}\left|\frac{Q(n+1)}{Q(n)}\right|=|k||z|$$
        To get the last equality, we used $\lim_{n\to\infty}\frac{Q(n+1)}{Q(n)}=1$, which holds because $Q(n+1)$ and $Q(n)$ have the same leading term and coefficient.
        
        From the above, we see that $\sum_{n=N+1}^\infty n! c_n z^n$ diverges whenever $|z|>1/|k|$. But the whole series $\sum_{n=0}^\infty n! c_n z^n$ is entire by hypothesis, so the same must true of the tail $\sum_{n=N+1}^\infty n!c_n z^n$. We have reached a contradiction, so $f$ must have infinitely many zeros.
    \end{proof}
    \begin{remark}
        Since the contradiction obtained in the proof of Theorem \ref{thm1.1} can be reached once we know $\sum_n n! c_n z^n$ converges at some $z_0\in\C$ with $|z_0|>1/|k|$, it is natural to ask whether we can determine the constant $k$ \textit{a priori}, that is, directly from the coefficients $c_n$. This would let us weaken the assumption that $\sum_n n! c_n z^n$ is entire to merely that it has radius of convergence $R>1/|k|$.
        
        Unfortunately, the utility of such a weakening is limited because Theorem \ref{thm1.1} can fail if $\sum_{n=0}^\infty n! c_n z^n$ is not entire. For any possible radius of convergence $0<R<\infty$, consider the series $\sum_{n=0}^\infty n! c_n z^n$ with the choice of Maclaurin coefficients
        $$c_n=\frac{1}{n! R^n}$$
        We can see that $c_n$ never terminates and the series
        $$\sum_{n=0}^\infty n! c_n z^n=\sum_{n=0}^\infty \frac{z^n}{R^n}$$
        has radius of convergence equal to $R$, so it is not entire. However, the corresponding $f$ is
        $$f(z)=\sum_{n=0}^\infty c_n z^n =\sum_{n=0}^\infty \frac{(z/R)^n}{n!}=e^{\frac{z}{R}}$$
        which does not have any zeros.
    \end{remark}
    \section{Proofs of Lemmas 1.2 and 1.3}
    In this section, we prove Lemmas \ref{lem1.2} and \ref{lem1.3}. We first prove \ref{lem1.2}, which is an easy consequence of the Hadamard factorization theorem and the generalized product rule
    $$(fg)^{(n)}(a)=\sum_{j=0}^n {\binom n j}f^{(j)}(a)g^{(n-j)}(a)$$
    for functions $f,g$ that are each $n$-times differentiable at $a$.
    \begin{proof}
        Since $f$ has order $1$ and $N$ zeros counting multiplicity, the Hadamard factorization of $f$ is $f(z)=e^{kz}P(z)$ for some constant $k\neq 0$, where $P(z)$ is a degree-$N$ polynomial. From the generalized Leibniz product rule,
        $$f^{(n)}(0)=\sum_{j=0}^n {\binom n j}k^{n-j}P^{(j)}(0)=k^n\sum_{j=0}^n {\binom n j}\frac{P^{(j)}(0)}{k^j}$$
        If $n>N$, we have $P^{(j)}(0)=0$ since $P$ is a polynomial of degree $N$, so
        $$f^{(n)}(0)=k^n\sum_{j=0}^N {\binom n j}\frac{P^{(j)}(0)}{k^j}\text{ for $n>N$}$$
        Write
        $$P(z)=\sum_{k=0}^N a_k z^k, a_k\in\C$$
        From the formula $a_j=\frac{P^{(j)}(0)}{j!}$ for the Maclaurin coefficients, we see that $P^{(j)}(0)=j!a_j$, so
        \begin{align*}
            f^{(n)}(0) &= k^n\sum_{j=0}^N {\binom n j}\frac{j!a_j}{k^j}\\
            &= k^n\sum_{j=0}^N \frac{n!}{(n-j)!}\cdot\frac{a_j}{k^j}\\
            &= k^n\sum_{j=0}^N \frac{a_j}{k^j}n(n-1)(n-2)\cdots(n-(j-1))
        \end{align*}
        The expressions
        $$\frac{a_j}{k^j}n(n-1)(n-2)\cdots(n-(j-1))$$
        are either $0$ or a polynomial (in $n$) of degree $j$, depending on whether $a_j=0$ or $a_j\neq 0$. Since $a_N\neq 0$ because $P$ has degree $N$, and since
        $$Q(n):=\sum_{j=0}^N \frac{a_j}{k^j}n(n-1)(n-2)\cdots(n-(j-1))$$
        is a finite sum of the above expressions, it must be a degree $N$ polynomial. This completes the proof of Lemma \ref{lem1.2}
    \end{proof}
    We now prove Lemma \ref{lem1.3}.
    \begin{proof}
        Suppose $\sum_{n=0}^\infty n!c_n z^n$ converges for $|z|<\delta$, where $\delta>0$. Evaluating at $z=\frac{\delta}{2}$ gives the convergent series $\sum_{n=0}^\infty n!c_n \left(\frac{\delta}{2}\right)^n$, so the terms $n!c_n \left(\frac{\delta}{2}\right)^n$ go to $0$ and are thus bounded. Suppose $M>0$ is such that $\left|n!c_n \left(\frac{\delta}{2}\right)^n\right|\leq M$. Then for all $z\in\C$,
        $$|c_nz^n|\leq \frac{M|z|^n\left(\frac{2}{\delta}\right)^n}{n!}=\frac{M\left(\frac{2|z|}{\delta}\right)^n}{n!}$$
        Summing the right side over $n\geq 0$ gives a series which converges to $Me^{2|z|/\delta}$, so $f$ is entire by comparison. We also see that
        \begin{equation}\label{ineq1.1}
            |f(z)|\leq \sum_{n=0}^\infty |c_n z^n|\leq M\sum_{n=0}^\infty\frac{\left(\frac{2|z|}{\delta}\right)^n}{n!}=Me^{2|z|/\delta}
        \end{equation}
        so $f$ has order at most $1$.
    \end{proof}
    \section{Applications of Theorem 1.1}\label{sec3}
   As promised, we now give alternative proofs that the Le Roy functions
   $$f_r(z)=\sum_{n=0}^\infty \frac{z^n}{(n!)^r}, r>1$$
   and Bessel functions
   $$J_\alpha(z)=\sum_{m=0}^\infty \frac{(-1)^m}{m!\Gamma(m+\alpha+1)}\left(\frac{z}{2}\right)^{2m+\alpha}$$
   have infinitely many zeros.
   For the Le Roy functions, we obtain a slightly stronger result: there are infinitely many zeros if $\on{Re}(r)>1$.
    \begin{theorem}
        If $\on{Re}(r)>1$, the function $f_r(z)=\sum_{n=0}^\infty \frac{z^n}{(n!)^r}$ is entire and has infinitely many zeros.
    \end{theorem}
    \begin{proof}
        By the ratio test, the series
        $$\sum_{n=0}^\infty n!\cdot \frac{z^n}{(n!)^r}=\sum_{n=0}^\infty \frac{z^n}{(n!)^{r-1}}$$
        is entire if $\on{Re}(r)>1$. The desired result follows from Theorem \ref{thm1.1}.
    \end{proof}
    We now prove that the Bessel functions have infinitely many zeros for any parameter $\alpha\in\R$.
    \begin{theorem}
        For each $\alpha\in\R$, the Bessel function
        $$J_\alpha(z)=\sum_{m=0}^\infty \frac{(-1)^m}{m!\Gamma(m+\alpha+1)}\left(\frac{z}{2}\right)^{2m+\alpha}$$
        has infinitely many zeros.
    \end{theorem}
    \begin{proof}
        Fix $\alpha\in\R$ and define
        $$c_m:=\frac{(-1)^m}{m!\Gamma(m+\alpha+1)2^{2m}}$$
        The series
        $$\sum_{m=0}^\infty m!c_m z^m=\sum_{m=0}^\infty \frac{(-1)^m}{\Gamma(m+\alpha+1)2^{2m}}z^m$$
        is entire by the ratio test, so $g(z):=\sum_{m=0}^\infty c_m z^m$ is entire and has infinitely many zeros by Theorem \ref{thm1.1}.
    
        Now,
        \begin{align*}
            J_\alpha(z) &= \sum_{m=0}^\infty \frac{(-1)^m}{m!\Gamma(m+\alpha+1)}\left(\frac{z}{2}\right)^{2m+\alpha}\\
            &= \left(\frac{z}{2}\right)^\alpha\sum_{m=0}^\infty \frac{(-1)^m}{m!\Gamma(m+\alpha+1)2^{2m}}(z^2)^m\\
            &= \left(\frac{z}{2}\right)^\alpha g(z^2)
        \end{align*}
        Let $a_1,a_2,\dots$ be the (infinitely many) zeros of $g(z)$. Then their principal square roots
        $$\sqrt{a_j}:=|a_j|^{\frac{1}{2}}e^{\frac{i\operatorname{Arg}(a_i)}{2}}$$
        are zeros of $g(z^2)$ and therefore also of $J_\alpha (z)$. It follows that $J_\alpha(z)$ also has infinitely many zeros.
    \end{proof}
    
	\bibliographystyle{plain} % Sets the bibliography style
	\bibliography{refs} % Bibliography entries are in the refs.bib file
	
	
	
	
\end{document}
